\documentclass[11pt,letterpaper]{article}
\usepackage[T1]{fontenc}
\usepackage{lmodern}
\usepackage[margin=1.05in]{geometry}
\usepackage{amsmath,amssymb,amsthm,booktabs,microtype,xurl}
\usepackage[hidelinks,pdftitle={Nonhalting and ZF-independence for a 209-state Turing machine},
  pdfauthor={James Chang}]{hyperref}

\newtheorem{theorem}{Theorem}[section]
\newtheorem{lemma}[theorem]{Lemma}
\newtheorem{definition}[theorem]{Definition}
\newcommand{\ZF}{\mathrm{ZF}}
\newcommand{\PA}{\mathrm{PA}}
\newcommand{\ISig}{\mathrm{I}\Sigma_1}
\newcommand{\N}{\mathbb{N}}
\newcommand{\Con}{\operatorname{Con}}
\newcommand{\Prov}{\operatorname{Prov}}
\newcommand{\Ic}{\mathsf{I}}
\newcommand{\Mc}{\mathsf{M}}
\newcommand{\Uc}{\mathsf{U}}
\newcommand{\Ec}{\mathsf{E}}
\newcommand{\Nc}{\mathsf{N}}
\newcommand{\Cc}{\mathsf{C}}
\newcommand{\mk}{\mathsf{Mark}}
\newcommand{\file}[1]{\texttt{\nolinkurl{#1}}}
\setlength{\emergencystretch}{2em}

\title{Nonhalting and ZF-independence\\for a 209-state Turing machine}
\author{James Chang}
\date{}

\begin{document}
\maketitle

\begin{abstract}
We give a two-symbol Turing machine with 209 working states that
searches for a contradiction in a small numeric calculus for set theory,
and a sentence of the language of set theory that expresses its
blank-tape halting in every model of ZF. If ZF is consistent, the machine
runs forever and ZF does not prove the negation of the sentence. If ZF is
arithmetically $\Sigma_1$-sound, ZF proves neither the sentence nor its
negation. The incompleteness step goes through an arithmetic theory whose
axioms are read off from proofs in the machine's calculus, so Foundation's
second incompleteness theorem for arithmetic applies directly. In Lean~4
we verify the calculus, the machine's simulation of its proof search, and
both unprovability results for the literal table.
\end{abstract}

\section{The result}\label{sec:result}

Let $M$ be the machine in Section~\ref{sec:machine}. Write $\ZF$ for
ordinary classical first-order Zermelo--Fraenkel set theory with full
Separation and Replacement and without choice. For an arithmetic sentence
$\theta$, write $\theta^\omega$ for its reading in the natural numbers of
a model of set theory.

\begin{definition}
$\ZF$ is \emph{arithmetically $\Sigma_1$-sound} if $\ZF\vdash\theta^\omega$
implies $\N\models\theta$ for every closed arithmetic $\Sigma_1$ sentence
$\theta$.
\end{definition}

\begin{theorem}\label{thm:main}
There is a closed sentence $H$ of the language of set theory such that,
in every model $V$ of $\ZF$, $V\models H$ exactly when $V$ contains an
internally finite halting history of $M$ from the blank tape. Moreover:
\begin{enumerate}
\item If $\ZF$ is consistent, $M$ does not halt and $\ZF\nvdash\neg H$.
\item If $\ZF$ is arithmetically $\Sigma_1$-sound, then $\ZF\nvdash H$
  and $\ZF\nvdash\neg H$.
\end{enumerate}
\end{theorem}

The histories may have nonstandard length. The meaning of $H$ and the
correctness of the compilation into the machine are conclusions of the
proof.

The machine searches for a refutation in a calculus $\Cc$ with
one set-theoretic axiom schema, the pointed gate
(Section~\ref{sec:calculus}). Soundness (Section~\ref{sec:sound}) shows
that a halting run gives a refutation of $\ZF$. The unprovability of
$\neg H$ is an instance of the following principle.

\begin{lemma}[Calibration]\label{lem:calib}
Let $T$ be an arithmetic theory with a $\Delta_1$ axiom set, and $\varphi$ a
set-theoretic sentence. Suppose that
\begin{enumerate}
\item $\ISig\subseteq T$;
\item if $\ZF$ is consistent, so is $T$;
\item $\ZF\vdash\neg\varphi\to\Con(T)^\omega$;
\item $\ZF\vdash\psi^\omega$ implies $T\vdash\psi$ for every arithmetic
  sentence $\psi$.
\end{enumerate}
If $\ZF$ is consistent, then $\ZF\nvdash\neg\varphi$.
\end{lemma}

\begin{proof}
If $\ZF\vdash\neg\varphi$, then (3) gives $\ZF\vdash\Con(T)^\omega$ and (4)
gives $T\vdash\Con(T)$. By (1), (2) and the second incompleteness
theorem~\cite{goedel} for $T$, this is impossible.
\end{proof}

Section~\ref{sec:theories} defines a theory $T$ from the machine's
calculus. Sections~\ref{sec:internal}--\ref{sec:absorb} prove the four
premises for $T$ and $\varphi=H$. Foundation~\cite{foundation} supplies
the second incompleteness theorem for arithmetic theories. No
derivability conditions for a set-theoretic provability predicate are
needed.

\section{The machine}\label{sec:machine}

The tape has alphabet $\{0,1\}$ and is infinite in both directions. It
starts with all cells zero. State $0$ is the start state, $0,\ldots,208$
are the working states, and $209$ is the halt label. In state $q$, a row
$(q,b,w,D,q')$ reads $b$, writes $w$, moves one cell in direction $D$, and
enters $q'$. The table has one row for each pair in
$\{0,\ldots,208\}\times\{0,1\}$. The file \file{data/zf209.tm} contains
all 418 rows and has SHA-256 digest
\begin{quote}\small
\nolinkurl{6b8acf3cf8523e2c2db2020f69c19f7ced2b1343e1b09621697515ed54b6eabd}.
\end{quote}
The first row is $(0,0,1,R,1)$, and $(62,1,1,R,209)$ is the only row that
enters the halt label. Appendix~\ref{app:table} lists the table.

As in~\cite{pa209}, a tape engine interprets a register program stored as
a binary tree. The program maintains a top formula, a coded stack and two
working registers. Its axiom, modus ponens and shift commands implement
the rules of $\Cc$, and its scheduler enumerates command words starting
from candidate $2$. Every $\Cc$-refutation has a command word that forces
halting from any retained store, so enumeration may keep the store
between candidates. Conversely, a halting command run supplies a
$\Cc$-refutation. From the blank tape the machine reaches the clean frame
of the scheduler's first site after $58{,}154$ steps.

\section{The calculus}\label{sec:calculus}

Let $p(a,b)=\binom{a+b+1}{2}+b$. The four constructors are
\[
 \Ic(a,b)=4p(a,b),\quad \Mc(a,b)=4p(a,b)+1,\quad
 \Uc(a,B)=4p(B,a)+2,\quad \Ec(a,b)=4p(a,b)+3,
\]
for implication, membership, universal quantification and equality.
A number used as an argument of $\Mc$, $\Ec$ or as the first argument of
$\Uc$ is a \emph{label}; $x_a$ is the variable with label $a$. Code $0$ is
truth and $1=\Mc(0,0)$ is falsity. Write $\Nc(B)=\Ic(B,1)$,
$B\wedge C=\Nc(\Ic(B,\Nc(C)))$, $B\leftrightarrow C$ for
$\Ic(B,C)\wedge\Ic(C,B)$, and $\exists x_a\,B$ for $\Nc(\Uc(a,\Nc(B)))$.
An arrow below denotes $\Ic$, and $x_a\in x_b$, $x_a=x_b$ denote
$\Mc(a,b)$, $\Ec(a,b)$.

A \emph{derivation} is a finite list in which each entry is an axiom, follows
from two earlier entries by modus ponens, or is $\Uc(x,A)$ for an earlier
entry $A$. The axioms are the six heads of Table~\ref{tab:heads}, the
restricted head
\[
 \rho_8(a,b,c,u)=\bigl((c\to u)\to a\bigr)\to\bigl((a\to c)\to(b\to c)\bigr)
 \qquad(u\le a),
\]
and the gate $F(a,b,c)$ for $c\ne a,b$.

\begin{table}[ht]
\centering\small
\begin{tabular}{@{}cl@{}}
\toprule
Head & Formula \\
\midrule
1 & $(\exists x_d\,\Uc(d,h))\to h$ \\
2 & $\Uc(d,\Uc(a,h))\to\Uc(a,\Uc(d,h))$ \\
3 & $x_h\in x_d\to(x_h=x_a\to x_a\in x_d)$ \\
4 & $x_h=x_d\to(x_a\in x_h\to x_a\in x_d)$ \\
5 & $\exists x_d\,(x_h=x_d)$ \\
6 & $\Uc(d,a)\to(\Uc(d,a\to h)\to\Uc(d,h))$ \\
\bottomrule
\end{tabular}
\caption{The six logical heads. The letters $a,d,h$ are numbers, read
as labels or as formula codes according to their position.}\label{tab:heads}
\end{table}

The gate uses the number $c$ both as a label and as a formula code. Write
$c$ for the formula, read with $x_b$ and $x_a$ as its first and second
arguments. Put
\begin{align*}
 \Delta&=x_a=x_b\to\Uc(c,x_a=x_b), &
 V&=x_a\in x_b\to\Uc(c,x_a\in x_b),
\end{align*}
and let $Q(a,b,c)=\Uc(c,G)$, where $G$ is the code of
\[
 \begin{aligned}
 &\Bigl[\forall x_a\Bigl(x_a\in x_c\to\neg\forall x_b\bigl(c\to\neg
  \forall x_b\bigl(\forall x_c(x_c\in x_b\to x_c\in x_a)\\
 &\qquad\qquad\qquad\qquad\quad{}\to
  \forall x_a(c\to\neg\forall x_a(c\to\neg\,x_a\in x_c))\bigr)\bigr)
 \Bigr)\Bigr]\\
 &\qquad{}\to\neg(c\to x_a\in x_c).
 \end{aligned}
\]
Then $F(a,b,c)=\Nc\bigl(\Delta\to(V\to Q)\bigr)$. Inside $G$ the binders
$\forall x_a$ and $\forall x_b$ rebind the arguments of $c$. Read
semantically, $\neg Q$ says that the relation $R(u,v)$ defined by $c$ has a
\emph{pointed hull} at $(x_b,x_a)$: a set $C$ with
\begin{equation}\label{eq:hull}
 R(x_b,x_a)\to x_a\in C,\qquad
 \forall z\in C\Bigl(\exists p\,R(p,z)\ \wedge\
 \forall u\subseteq z\bigl(\exists v\,R(u,v)\to\exists w\in C\,R(u,w)\bigr)\Bigr).
\end{equation}

\section{Soundness}\label{sec:sound}

\begin{lemma}[Definable pointed hulls]\label{lem:hull}
For every definable relation $R$ with parameters and all $u_0,v_0$, $\ZF$
proves that some set $C$ satisfies~\eqref{eq:hull} at $(u_0,v_0)$.
\end{lemma}

\begin{proof}
If $R(u,\cdot)$ is nonempty, let $\alpha_u$ be the least stage of the
cumulative hierarchy meeting it and $S(u)=\{v\in V_{\alpha_u}:R(u,v)\}$;
otherwise $S(u)=\varnothing$. Start with $C_0=\{v\in\{v_0\}:R(u_0,v)\}$ and
put $C_{n+1}=C_n\cup\bigcup\{S(u):\exists z\in C_n\,(u\subseteq z)\}$.
Power Set and Replacement make each stage a set, and $C=\bigcup_nC_n$ is
as required. No choice is used.
\end{proof}

\begin{lemma}\label{lem:sound}
If $M$ halts from the blank tape, $\ZF$ is inconsistent.
\end{lemma}

\begin{proof}
A halting run gives a $\Cc$-refutation (Section~\ref{sec:machine}). In a
model of $\ZF$, interpret membership and equality as themselves, with
label $0$ assigned the empty set so that code $1$ is false. The heads and
$\rho_8$ are valid. Lemma~\ref{lem:hull} gives $\neg Q$; the guard
$c\ne a,b$ and the fact that every label of a code is below the code make
$\Delta$ and $V$ valid. The rules preserve validity, so a refutation
contradicts the existence of a model. The reverse translation into a
fixed presentation of $\ZF$ makes this an actual finite $\ZF$ refutation.
\end{proof}

\section{The theories}\label{sec:theories}

\paragraph{The translation $\chi$.}
A set-theoretic formula becomes a $\Cc$ code by a translation $\chi$ that
keeps track of labels. Free variable $i$ gets the label $Y(m+i)$; a
quantifier at depth $d$ binds the label $X(f,d)$, where $f$ names a label
family; three reserved labels $R_0,R_1,R_2$ are used only inside atoms.
These are distinct numbers given by pairing, and none is $0$. Equality
becomes the defined equality
\[
 x\approx y\ =\ \forall r_0(r_0\in x\leftrightarrow r_0\in y)\wedge
                \forall r_0(x\in r_0\leftrightarrow y\in r_0),
\]
membership becomes $\exists r_1\exists r_2(x\approx r_1\wedge y\approx r_2
\wedge r_1\in r_2)$, and connectives and quantifiers become their $\Cc$
codes. The translation computes the pair $(\chi\varphi,\chi\neg\varphi)$
at once.

\paragraph{The theories $U$ and $T$.}
Let $U$ be the set-theoretic theory with axioms $\theta\wedge\top^n$, where
$n$ codes a $\Cc$-derivation containing $\chi\theta$. Let $T$ be the
arithmetic theory with axioms $\sigma\wedge\top^n$, where $n$ codes a
$\Cc$-derivation containing $\chi(\kappa\sigma)$ and
\[
 \kappa\sigma\ =\ \forall e\,\bigl(\mk(e)\to J(e,\sigma)\bigr).
\]
Here $J(e,\cdot)$ reads arithmetic formulas in set theory through seven
fixed set-theoretic templates with a marker $e$: a domain, equality, order,
zero, one, addition and multiplication. $\mk(e)$ says that $e$ has no
members. The witness $n$ lies below the axiom, and checking a derivation
is $\Delta_1$, so both axiom sets are $\Delta_1$ (Craig's trick).

The templates are chosen as follows. At a marker $e$ with no members, a
structure satisfying the gate contains a model $N_e$ of $\ZF$, defined in
three layers: indiscernibility classes of the base structure, marked
objects modulo marked bisimilarity with marker $e$, and the well-founded
part. The gate supplies the Separation, Replacement and closure
constructions each layer needs. Each object is projected to the natural number of $N_e$
it denotes, or to zero. The templates relate projections through the
formulas that read arithmetic in $\omega^{N_e}$. At any other $e$ every
template holds.

Two properties of $U$ follow from its definition. If $U\vdash\theta$, then
$\Cc$ proves $\chi\theta$ (Section~\ref{sec:internal}, applied in $\N$).
If $\Cc$ proves $\chi\theta$, then $\theta\wedge\top^n$ is an axiom of $U$.
Similarly, a $\Cc$-proof of $\chi(\kappa\sigma)$ puts $\sigma\wedge\top^n$
into $T$.

\section{The internal premise}\label{sec:internal}

Let $H_{\mathrm{phys}}$ be the arithmetic sentence asserting the existence
of a coded blank-tape history for the literal table, with an initial blank
configuration, valid adjacent steps, and a terminal configuration. Write
$H=H_{\mathrm{phys}}^\omega$. Lean also gives a $\Sigma_1$ sentence
$H_\Sigma$ for the same halting predicate and proves
$\ZF\vdash H\to H_\Sigma^\omega$.

\begin{lemma}\label{lem:internal}
$\PA\vdash\neg H_{\mathrm{phys}}\to\Con(T)$. Hence premise~(3) of
Lemma~\ref{lem:calib} holds
for $\varphi=H$.
\end{lemma}

\begin{proof}
Work in a model of $\PA$ and suppose it contains a coded $T$-refutation,
possibly nonstandard. We transport it in two stages.

\emph{Stage 1.} The translation $J(\varepsilon,\cdot)$, with a free marker
$\varepsilon$, sends the refutation to a $U$-derivation of $\neg\mk(\varepsilon)$.
The axiom $\sigma\wedge\top^n$ of $T$ comes from the axiom
$\kappa\sigma\wedge\top^n$ of $U$ by instantiating $e$. The quantifier steps
need a substitution lemma for $J$, proved by induction on formulas from
seven template facts: the domain is everything, equality is an equivalence,
the operations are total and functional, and every template respects
equality. Section~\ref{sec:absorb} shows that $U$ proves these facts and
$\exists e\,\mk(e)$; together with $\forall e\,\neg\mk(e)$ this refutes $U$.

\emph{Stage 2.} Every internal derivation of $U$ becomes a $\Cc$-proof of
the translation of its conclusion. Each logical rule has a fixed template
in $\Cc$. Eigenvariables take free labels from a window that the derivation
code bounds, and the label discipline of $\chi$ supplies the vacuity and
substitution side conditions of the quantifier rules. Axioms of $U$ come
with their $\Cc$-proofs.

A $\Cc$-proof of $1$ halts the machine: the scheduler reaches its command
word, and the tape simulation, including the startup, gives a coded halting
run of the literal table. All steps hold in arbitrary models of $\PA$,
so first-order completeness gives the displayed implication. Interpreting
it in $\omega$ gives $\ZF\vdash\neg H\to\Con(T)^\omega$.
\end{proof}

\section{The gate inside \texorpdfstring{$U$}{U}}\label{sec:gate}

The template facts and absorption are proved semantically, in models of
$U$. What a model $W$ of $U$ needs is the gate: every definable relation
on its indiscernibility quotient has pointed hulls. Then $W$ has an object
with no members, and each such object $e$ gives the model $N_e$ of $\ZF$.

For a set-theoretic formula $\varphi(D,H,\bar p)$ with $k$ parameters, let
$\gamma_\varphi$ say: for all $\bar p$ and all $u,v$, some $C$ is a pointed
hull of $R=\varphi(\cdot,\cdot,\bar p)$ at $(u,v)$, written with the exact
connectives and binders of $\neg Q$.

\begin{lemma}\label{lem:gatesent}
$\Cc$ proves $\chi\gamma_\varphi$. Hence $U\vdash\gamma_\varphi$.
\end{lemma}

\begin{proof}
Take labels $a=X(0,k+1)$ and $b=X(0,k)$, the labels $\chi$ gives the
binders of $u$ and $v$. Let the relation code $c$ be $\chi\varphi$ in the
separate family $X(1,\cdot)$, conjoined with a closed theorem of $\Cc$
whose binders are $R_0,R_1,R_2$ and $X(0,j)$ for $j<k+10$. Then $c$
exceeds every label the argument must avoid, and the gate instance
$\neg Q(a,b,c)$ is available.

Three steps turn $\neg Q(a,b,c)$ into $\chi\gamma_\varphi$.
First, $\Cc$ proves $x\in y\leftrightarrow\chi(x\in y)$ for distinct labels:
the witnesses come from head~5 and the direct rewrites, and each conjunct
of $\approx$ is used at a label introduced by head~5, with the backward
rewrites. Second, a renaming theorem relabels every binder of $\neg Q$
with the label $\chi$ uses at that depth. It is proved once for a small
syntax of gate-shaped formulas: under a substitution and vacuity
interface for the relation code, and under conditions that no binder
captures a variable used beneath it, the two labellings are equivalent in
$\Cc$. Third, at each occurrence of the relation, $\chi$ of the substituted
formula is $\chi\varphi$ with relabelled arguments at the current depth in
family $0$. Renaming between families and depths gives equivalence with
$c$. The substitution and vacuity conditions for $\chi\varphi$ need only the
labels the translation actually binds, and these lie below $c$.
\end{proof}

\begin{lemma}\label{lem:Uhull}
In every model $W$ of $U$, every relation definable with parameters on the
indiscernibility quotient of $W$ has pointed hulls, and $W$ has an object
with no members.
\end{lemma}

\begin{proof}
Lemma~\ref{lem:gatesent} gives hulls in $W$ for relations defined with
parameters. A relation on the quotient defined by $\varphi$ is defined in
$W$ by $\varphi$ with equality replaced by indiscernibility, so its hull in
$W$ projects to a hull in the quotient. No extensionality is needed in $W$.
The empty relation has a hull, which has no members.
\end{proof}

\section{Absorption and consistency}\label{sec:absorb}

\begin{lemma}\label{lem:absorb}
Premises (1), (2) and (4) of Lemma~\ref{lem:calib} hold for $T$. Moreover
$U$ proves the template facts and $\exists e\,\mk(e)$.
\end{lemma}

\begin{proof}
Let $W$ be a model of $U$. By Lemma~\ref{lem:Uhull} and
Section~\ref{sec:theories}, every marker $e$ of $W$ gives a model $N_e$ of $\ZF$, and the templates at
$e$ read the arithmetic of $\omega^{N_e}$ through the projection. The
template facts hold at such $e$ and trivially at other $e$. By
completeness, $U$ proves them, and Lemma~\ref{lem:Uhull} gives
$\exists e\,\mk(e)$.

(4) If $\ZF\vdash\psi^\omega$, then $\psi$ holds in $\omega^{N_e}$ for every
marker $e$ of every model of $U$, so $W\models\kappa\psi$. Hence
$U\vdash\kappa\psi$, $\Cc$ proves $\chi(\kappa\psi)$, and $T\vdash\psi$.
(1) Every axiom of $\PA$ is proved by $\ZF$ in $\omega$, so (4) gives
$\PA\subseteq T$ and $\ISig\subseteq T$.

(2) Let $V$ be a model of $\ZF$. Every theorem of $\Cc$ holds in $V$ with
label $0$ assigned the empty set (Lemma~\ref{lem:sound}), and the reading
of $\chi\theta$ in $V$ is the reading of $\theta$. So $V\models\kappa\sigma$
for every axiom $\sigma\wedge\top^n$ of $T$. At the empty set $e$ of $V$,
the templates read $\omega^{N_e}$, which is therefore a model of $T$.
\end{proof}

\begin{proof}[Proof of Theorem~\ref{thm:main}]
The interpretation of $H_{\mathrm{phys}}$ in $\omega$ gives the internally
finite history in the statement of the theorem. Under consistency,
Lemma~\ref{lem:sound} gives nonhalting. Lemmas~\ref{lem:internal}
and~\ref{lem:absorb} give the four premises of Lemma~\ref{lem:calib} for
$T$ and $H$, so $\ZF\nvdash\neg H$.

Under arithmetic $\Sigma_1$-soundness, $\ZF$ is consistent, so the first
part applies. If $\ZF\vdash H$, the implication to $H_\Sigma^\omega$
gives $\ZF\vdash H_\Sigma^\omega$. By soundness $\N\models H_\Sigma$,
so $M$ halts, a contradiction.
\end{proof}

\section{Lean verification}\label{sec:status}

The standalone \file{ProofSources/Challenge.lean} uses only \file{Std}.
It defines the proof system of $\ZF$, its axioms, the literal table and the
model-internal meaning of halting, and states the claim as
\file{PalomarZF209.exact209Claim}. \file{ProofSources/Solution.lean} proves
it from \file{Halting/Final.lean}; its one statement placeholder lies
outside the Solution's imports. Paths below are relative to
\file{ProofSources/}.

The calibration principle is \file{Independence/Godel.lean}, and
\file{Independence/Top.lean} derives the compared claim from it.
\file{Halting/Theories.lean} and \file{Halting/Craig.lean} define $U$ and
$T$. \file{Internal/Arith/} contains stage~1 and
\file{Internal/Tait/Derivation.lean} stage~2; \file{Halting/Premise.lean}
combines them with \file{Halting/Prov.lean}, which turns a $\Cc$-proof of
$1$ into halting of the literal table. \file{Halting/GateAtom.lean},
\file{Halting/GateAlpha.lean}, \file{Halting/GateChi.lean},
\file{Halting/GateFound.lean} and \file{Halting/GateC.lean} prove
Lemma~\ref{lem:gatesent}; \file{Halting/GateHull.lean} proves
Lemma~\ref{lem:Uhull}. \file{GateModel/} builds $N_e$, and
\file{Halting/Absorb.lean} and \file{Halting/Consistent.lean} prove
Lemma~\ref{lem:absorb}. \file{ScottHull.lean} proves
Lemma~\ref{lem:hull}, and \file{Literal209FoundationNonhalting.lean} the
nonhalting theorem. \file{M209/} contains the startup and the simulation
of the register program.

The package has 1,225 Lean files with 156,142 lines, including 161
vendored Foundation modules. Each Lean process is capped at 1,536 MiB.
A clean build of the package, with Mathlib's cached build, took
14.4 minutes on a ten-core machine, and the largest process used 1.7 GB.
The axiom audit reports only \file{propext}, \file{Classical.choice} and
\file{Quot.sound}. The Solution uses no \file{sorry},
\file{native_decide} or added axioms. \file{docs/reproduction.md} gives the
build, audit and Comparator commands.

The formalization uses Foundation~\cite{foundation} and a tape engine
ported from Elliot Glazer's K230 Rocq development, with components adapted
from Jeff Ketchersid's Busy Beaver developments. \file{docs/provenance.md}
records the revisions, licences and source notices.

\begin{thebibliography}{9}

\bibitem{pa209}
J. Chang.
Nonhalting and PA-independence for a 209-state Turing machine.
\href{https://github.com/jamesyc/BusyBeaver/tree/69d98241a721ead20b4f0b17df9062950b40fbfe}{Public PA209 snapshot}.

\bibitem{foundation}
FormalizedFormalLogic.
Foundation.
\url{https://github.com/FormalizedFormalLogic/Foundation}.

\bibitem{goedel}
K. G\"odel.
\"Uber formal unentscheidbare S\"atze der Principia Mathematica und
verwandter Systeme I.
\emph{Monatshefte f\"ur Mathematik und Physik} 38 (1931), 173--198.
\url{https://doi.org/10.1007/BF01700692}.

\end{thebibliography}

\appendix
\section{The literal table}\label{app:table}

Each entry below is ``write, move, successor''. The column heading is the
read symbol.

\begin{center}
\scriptsize
\setlength{\tabcolsep}{4pt}
\begin{tabular}{rll@{\qquad}rll@{\qquad}rll@{\qquad}rll@{\qquad}rll}
\toprule
$q$ & $0$ & $1$ & $q$ & $0$ & $1$ & $q$ & $0$ & $1$ & $q$ & $0$ & $1$ & $q$ & $0$ & $1$ \\
\midrule
0 & \texttt{1R1} & \texttt{1R60} & 42 & \texttt{0R54} & \texttt{1R55} & 84 & \texttt{1L98} & \texttt{1L84} & 126 & \texttt{0R43} & \texttt{1R89} & 168 & \texttt{0R142} & \texttt{1R160} \\
1 & \texttt{1R81} & \texttt{1R1} & 43 & \texttt{0R23} & \texttt{0R51} & 85 & \texttt{0R18} & \texttt{1R99} & 127 & \texttt{0L5} & \texttt{0L127} & 169 & \texttt{0R167} & \texttt{1R168} \\
2 & \texttt{0L2} & \texttt{1L3} & 44 & \texttt{0R57} & \texttt{1R58} & 86 & \texttt{0R57} & \texttt{1R57} & 128 & \texttt{0R125} & \texttt{1R58} & 170 & \texttt{0R130} & \texttt{1R169} \\
3 & \texttt{0L4} & \texttt{1L3} & 45 & \texttt{0R59} & \texttt{0L15} & 87 & \texttt{0R115} & \texttt{1R114} & 129 & \texttt{0R11} & \texttt{1R36} & 171 & \texttt{0R63} & \texttt{1R170} \\
4 & \texttt{0L6} & \texttt{1L3} & 46 & \texttt{0R24} & \texttt{1R46} & 88 & \texttt{0L25} & \texttt{1R208} & 130 & \texttt{0R18} & \texttt{1R19} & 172 & \texttt{0R57} & \texttt{1R137} \\
5 & \texttt{1R7} & \texttt{0L5} & 47 & \texttt{0L34} & \texttt{0R60} & 89 & \texttt{0R94} & \texttt{0L127} & 131 & \texttt{0R77} & \texttt{1R86} & 173 & \texttt{0R147} & \texttt{1R162} \\
6 & \texttt{0L9} & \texttt{1L3} & 48 & \texttt{0R15} & \texttt{1R48} & 90 & \texttt{0R101} & \texttt{1R102} & 132 & \texttt{0R85} & \texttt{1R131} & 174 & \texttt{0R172} & \texttt{1R173} \\
7 & \texttt{0R7} & \texttt{0R10} & 49 & \texttt{0R62} & \texttt{1R63} & 91 & \texttt{0R66} & \texttt{1R55} & 133 & \texttt{0R130} & \texttt{1R132} & 175 & \texttt{0R171} & \texttt{1R174} \\
8 & \texttt{0R13} & \texttt{1R12} & 50 & \texttt{0L38} & \texttt{1L50} & 92 & \texttt{0R54} & \texttt{1R103} & 134 & \texttt{0L159} & \texttt{1R77} & 176 & \texttt{0R175} & \texttt{1R133} \\
9 & \texttt{0L14} & \texttt{1L3} & 51 & \texttt{0R21} & \texttt{1R51} & 93 & \texttt{0R20} & \texttt{1R104} & 135 & \texttt{0L169} & \texttt{1R134} & 177 & \texttt{0R133} & \texttt{1R176} \\
10 & \texttt{0L15} & \texttt{0R10} & 52 & \texttt{0L11} & \texttt{0R48} & 94 & \texttt{0R20} & \texttt{0L10} & 136 & \texttt{0R18} & \texttt{1R116} & 178 & \texttt{0R160} & \texttt{1R57} \\
11 & \texttt{0L25} & \texttt{1R17} & 53 & \texttt{0R57} & \texttt{0L15} & 95 & \texttt{0R58} & \texttt{1R40} & 137 & \texttt{0R136} & \texttt{1R134} & 179 & \texttt{0R130} & \texttt{1R178} \\
12 & \texttt{0L78} & \texttt{1L30} & 54 & \texttt{0R64} & \texttt{1R65} & 96 & \texttt{0R105} & \texttt{1R8} & 138 & \texttt{0R134} & \texttt{1R133} & 180 & \texttt{0R179} & \texttt{1R161} \\
13 & \texttt{0L11} & \texttt{0R21} & 55 & \texttt{0R58} & \texttt{1R66} & 97 & \texttt{0R90} & \texttt{1R106} & 139 & \texttt{0R57} & \texttt{1R77} & 181 & \texttt{0R148} & \texttt{1R180} \\
14 & \texttt{0L22} & \texttt{1L3} & 56 & \texttt{0R51} & \texttt{1R56} & 98 & \texttt{0L127} & \texttt{0L84} & 140 & \texttt{0R18} & \texttt{1R139} & 182 & \texttt{0R151} & \texttt{1R181} \\
15 & \texttt{1L2} & \texttt{0L2} & 57 & \texttt{0R47} & \texttt{0R24} & 99 & \texttt{0R107} & \texttt{1R93} & 141 & \texttt{0L197} & \texttt{1R140} & 183 & \texttt{0R153} & \texttt{1R182} \\
16 & \texttt{0R23} & \texttt{0R31} & 58 & \texttt{0L27} & \texttt{1L67} & 100 & \texttt{0L108} & \texttt{0L2} & 142 & \texttt{0R140} & \texttt{1R141} & 184 & \texttt{0R139} & \texttt{1R145} \\
17 & \texttt{0R83} & \texttt{0L94} & 59 & \texttt{0L101} & \texttt{1L29} & 101 & \texttt{0R37} & \texttt{1R67} & 143 & \texttt{0R138} & \texttt{1R135} & 185 & \texttt{0R137} & \texttt{1R184} \\
18 & \texttt{0R27} & \texttt{1R28} & 60 & \texttt{0R0} & \texttt{1R60} & 102 & \texttt{0R58} & \texttt{1R109} & 144 & \texttt{0R137} & \texttt{1R143} & 186 & \texttt{0R136} & \texttt{1R179} \\
19 & \texttt{0R29} & \texttt{1R30} & 61 & \texttt{1R123} & \texttt{1L50} & 103 & \texttt{0R58} & \texttt{1R110} & 145 & \texttt{0R136} & \texttt{1R131} & 187 & \texttt{0R140} & \texttt{1R77} \\
20 & \texttt{0R23} & \texttt{0R56} & 62 & \texttt{0R69} & \texttt{1R209} & 104 & \texttt{0R65} & \texttt{0L15} & 146 & \texttt{0R130} & \texttt{1R145} & 188 & \texttt{0R187} & \texttt{1R149} \\
21 & \texttt{0R32} & \texttt{1R21} & 63 & \texttt{0L124} & \texttt{1R30} & 105 & \texttt{0R13} & \texttt{1R109} & 147 & \texttt{0R132} & \texttt{1R133} & 189 & \texttt{0R186} & \texttt{1R188} \\
22 & \texttt{0L33} & \texttt{1L3} & 64 & \texttt{0R39} & \texttt{1R59} & 106 & \texttt{0R20} & \texttt{1R111} & 148 & \texttt{0R145} & \texttt{1R77} & 190 & \texttt{0R136} & \texttt{1R189} \\
23 & \texttt{0L34} & \texttt{0R1} & 65 & \texttt{0L78} & \texttt{1R70} & 107 & \texttt{0R112} & \texttt{1L19} & 149 & \texttt{0L147} & \texttt{1L184} & 191 & \texttt{0R190} & \texttt{1R188} \\
24 & \texttt{0R35} & \texttt{1R24} & 66 & \texttt{0R52} & \texttt{1L71} & 108 & \texttt{0R29} & \texttt{1R64} & 150 & \texttt{0R149} & \texttt{1R57} & 192 & \texttt{0R185} & \texttt{1R191} \\
25 & \texttt{0R120} & \texttt{1R40} & 67 & \texttt{0R72} & \texttt{1R73} & 109 & \texttt{0R37} & \texttt{1L107} & 151 & \texttt{0R130} & \texttt{1R150} & 193 & \texttt{0R183} & \texttt{1R192} \\
26 & \texttt{0R37} & \texttt{0L38} & 68 & \texttt{0R68} & \texttt{1R74} & 110 & \texttt{0R119} & \texttt{0L15} & 152 & \texttt{0R148} & \texttt{1R151} & 194 & \texttt{0L197} & \texttt{1R187} \\
27 & \texttt{0R39} & \texttt{1R40} & 69 & \texttt{0R75} & \texttt{1R76} & 111 & \texttt{0R113} & \texttt{0L15} & 153 & \texttt{0L174} & \texttt{1R152} & 195 & \texttt{0R77} & \texttt{1R194} \\
28 & \texttt{0R41} & \texttt{1R42} & 70 & \texttt{0R78} & \texttt{1L109} & 112 & \texttt{0R23} & \texttt{0R46} & 154 & \texttt{0R146} & \texttt{1R153} & 196 & \texttt{0R132} & \texttt{1R188} \\
29 & \texttt{0R43} & \texttt{1R44} & 71 & \texttt{0L97} & \texttt{1R53} & 113 & \texttt{0R47} & \texttt{0R31} & 155 & \texttt{0R144} & \texttt{1R154} & 197 & \texttt{0R195} & \texttt{1R196} \\
30 & \texttt{0R20} & \texttt{1R45} & 72 & \texttt{0R47} & \texttt{0R32} & 114 & \texttt{0R69} & \texttt{1R85} & 156 & \texttt{0R135} & \texttt{1R155} & 198 & \texttt{0R193} & \texttt{1R197} \\
31 & \texttt{0R46} & \texttt{1R31} & 73 & \texttt{0R80} & \texttt{0L40} & 115 & \texttt{0R124} & \texttt{1R27} & 157 & \texttt{0R124} & \texttt{1R156} & 199 & \texttt{0R194} & \texttt{1R148} \\
32 & \texttt{0R31} & \texttt{1R32} & 74 & \texttt{0R49} & \texttt{1R128} & 116 & \texttt{0R120} & \texttt{1R106} & 158 & \texttt{0R77} & \texttt{1R157} & 200 & \texttt{0R141} & \texttt{1R120} \\
33 & \texttt{0L68} & \texttt{1L3} & 75 & \texttt{0R126} & \texttt{1R20} & 117 & \texttt{0R93} & \texttt{1R85} & 159 & \texttt{0R158} & \texttt{1R146} & 201 & \texttt{0R121} & \texttt{0L127} \\
34 & \texttt{1R34} & \texttt{1L2} & 76 & \texttt{0R100} & \texttt{1R89} & 118 & \texttt{0R82} & \texttt{1R97} & 160 & \texttt{0R136} & \texttt{1L152} & 202 & \texttt{0R122} & \texttt{1R201} \\
35 & \texttt{0R48} & \texttt{1R35} & 77 & \texttt{0R91} & \texttt{1R92} & 119 & \texttt{0R86} & \texttt{1R86} & 161 & \texttt{0L184} & \texttt{1R160} & 203 & \texttt{0R121} & \texttt{1R202} \\
36 & \texttt{0R27} & \texttt{1R19} & 78 & \texttt{0R72} & \texttt{1R71} & 120 & \texttt{0R16} & \texttt{1L108} & 162 & \texttt{0R138} & \texttt{1R57} & 204 & \texttt{0R200} & \texttt{1R203} \\
37 & \texttt{0L11} & \texttt{0R35} & 79 & \texttt{0R90} & \texttt{1R117} & 121 & \texttt{0R69} & \texttt{1R100} & 163 & \texttt{0R161} & \texttt{1R162} & 205 & \texttt{0R199} & \texttt{1R204} \\
38 & \texttt{0L5} & \texttt{1L50} & 80 & \texttt{0R94} & \texttt{1R95} & 122 & \texttt{0R95} & \texttt{1R107} & 164 & \texttt{0R163} & \texttt{1R137} & 206 & \texttt{0R186} & \texttt{1R205} \\
39 & \texttt{0R23} & \texttt{0R24} & 81 & \texttt{0R61} & \texttt{0R1} & 123 & \texttt{1L38} & \texttt{1L98} & 165 & \texttt{0R164} & \texttt{1R163} & 207 & \texttt{0R198} & \texttt{1R206} \\
40 & \texttt{0R12} & \texttt{0L15} & 82 & \texttt{0R26} & \texttt{1R96} & 124 & \texttt{0R79} & \texttt{1R77} & 166 & \texttt{0R159} & \texttt{1R165} & 208 & \texttt{0R177} & \texttt{1R207} \\
41 & \texttt{0R52} & \texttt{1R53} & 83 & \texttt{0R88} & \texttt{1R87} & 125 & \texttt{0R118} & \texttt{1R129} & 167 & \texttt{0R85} & \texttt{1R166} &  &  &  \\
\bottomrule
\end{tabular}
\end{center}


\end{document}
